\begin{proof}
\begin{enumerate}[label=\arabic*.]
\item We remove the existential free-lift condition by an explicit
construction. Let $d>1$ be the order of $M$ and choose a primitive
$u\in\ker(M-I)\cap\mathbb Z^6$. Such a vector is obtained from a
nonzero rational kernel vector by clearing denominators and dividing
by the greatest common divisor of its integer coordinates; both
operations retain the equation $(M-I)u=0$. Put
$g(x)=Mx+u/d^2$ and $S_d=\sum_{a=0}^{d-1}M^a$.
Then $g^j(x)=M^jx+ju/d^2$, $S_dM^j=S_d$ and $S_du=du$.
A fixed point of $g^j$ gives
$(M^j-I)x+ju/d^2=\lambda\in\mathbb Z^6$.
Applying $S_d$ gives $ju/d=S_d\lambda\in\mathbb Z^6$.
A Bezout row pairing to one with $u$ therefore gives $d\mid j$.
Writing $j=dr$, the original equation becomes $ru/d=\lambda$,
so the same row gives $d\mid r$. Thus only multiples of $d^2$
can have fixed points. Since $g^{d^2}$ is translation by the
integer vector $u$, its torus action is the identity. This proves
both its actual order $d^2$ and freeness of every nonidentity power.
For $M=I$, translation by $u/2$ for primitive $u$ has actual order
two and is free. Conversely, when $I-M$ is invertible over
$\mathbb R$, every affine map $Mx+a$ fixes the torus class of
$(I-M)^{-1}a$. The free-lift test is consequently exactly
$\ker(M-I)\ne0$.

\item The finite linear orders here are the already proved list of
\cref{maps:src:F.2.2}:
\[
 d\in\{1,2,3,4,5,6,8,10,12\}.
 \tag{F.CN1}\label{maps:eq:F.CN1}
\]
Taking these finitely many powers, checking the fixed kernel,
the product relation and the cusp Smith factors therefore gives
an exact admissibility test. The actual affine order is the
separate retained gluing datum; the construction in item1 gives
its exact relation to the chosen finite lift.

\item For an admissible full tuple $T=(M_1,\ldots,M_k,M_0)$ let
$H(T)$ be its maximum absolute entry. There is at most one identity
finite slot in a three-point tuple: two such identities would make
$M_0$ finite order, whereas $M_0=I+N$, $N^2=0$ and
$\operatorname{rank}N=3$ give $(I+N)^d=I+dN\ne I$ for every
positive integer $d$. At a fixed identity slot, delete that matrix
and reinsert $I$ at the same slot after classifying the remaining
ordered tuple. These two maps are inverse, and both commute with
$T\mapsto (QM_1Q^{-1},\ldots,QM_0Q^{-1})$ for the same
$Q\in\mathrm{GL}(6,\mathbb Z)$. They therefore retain the actual
returned simultaneous conjugator. The remaining invertible integral
matrices have height at least one, so inserting $I$ leaves height
unchanged. Every entry in the inserted block is constant; the first
unequal row-major entry of two tuples is therefore the same before
and after insertion. This proves preservation of the specified
lexicographic comparison as well. The actual affine order and
numerator at that slot are retained as gluing data and are transported
by this same $Q$. In the two-involution sector use the canonical
representative of \ref{maps:src:F.2.38}, with its own specified
convention. In all remaining sectors, order full integer tuples
first by $H$, then lexicographically by their row-major entries in
the displayed marked order, and take the first admissible member
of the simultaneous conjugacy class.

\item To compute the representative, enumerate heights
$1,\ldots,H(T)$ and all integer entries in the stated order.
There are at most
\[
 (2H(T)+1)^{36(k+1)}
 \tag{F.CN2}\label{maps:eq:F.CN2}
\]
integer lists in this box. Retain the invertible tuples satisfying
the ordered product relation and the stated maximal square-zero
primitive cusp condition, which are finite integer tests, and
apply the decision-and-witness algorithm of
\cref{maps:src:F.2.36} to each candidate and the supplied tuple.
Its general rank-three square-zero input permits precisely this
call. A positive answer transfers all finite-order, orientation
and admissibility conditions through its actual integral
conjugator. The supplied tuple itself occurs, with the identity
as a possible witness, so a positive answer is reached in
finitely many tests. The first one is the required representative.
Record the returned convention as
$P_T^{-1}T_iP_T=\operatorname{Can}(T)_i$.
The displayed candidate bound makes no assumption about entries
of $P_T$; the terminating witness construction is the stated
input from \cref{maps:src:F.2.36}.

\item Every simultaneous orbit has a least element in this order.
Its height is at most the height of each supplied member, so
\eqref{maps:eq:F.CN2} includes that global least element.
Equivalent inputs have exactly the same ordered orbit and
therefore the same representative. Conversely, if $T,T'$ have
the same representative, the actual matrix
$Q=P_TP_{T'}^{-1}$ satisfies $Q^{-1}T_iQ=T'_i$ for every slot.
The separate involution convention has the same property by
\cref{maps:src:F.2.38}, with the product of its actual returned
witnesses giving the corresponding map. Conjugation preserves
the ordered finite orders, so the involution sector is preserved
by those same maps. This proves the classification assertion
and its finite stopping rule.

\item Finally, \ref{maps:src:F.2.16} already gives an unbounded admissible integral
family. In its homology convention direct symbolic
multiplication gives
\[
             \operatorname{tr}(M_1M_2^3)
                       =-v^4+v^3-v^2+2.                 \tag{F.CN3}\label{maps:eq:F.CN3}
\]
This is unbounded even on the previously constructed free
affine subfamily $v\in20\mathbb Z$. It is invariant under
simultaneous conjugacy. A fixed finite entry box contains
only finitely many tuples and hence only finitely many such
traces; it cannot contain a representative of every class.
The input-dependent nature of \eqref{maps:eq:F.CN2} is therefore necessary
for this particular exhaustive procedure.

\end{enumerate}
\end{proof}
